\cvsection{Expériences professionnelles}

\begin{cventries}
  \cventry
    {Stagiaire ingénieur de recherche en IA}
    {Inria}
    {Bordeaux, France}
    {Mars 2026 -- Septembre 2026}
    {
      \begin{cvitems}
        \item {Conception et développement en autonomie, avec le suivi de mes encadrants, d’un agent IA multi-étapes répondant au besoin défini par mon tuteur pour aider les chercheurs à trouver des articles scientifiques.}
        \item {Intégration de l’API OpenAI avec création et orchestration des outils de l’agent.}
        \item {Développement de services et d’API REST en Python/FastAPI avec MongoDB, ainsi que d’une interface React/Next.js permettant aux chercheurs d’interagir avec l’agent.}
        \item {Mise en œuvre d’un pipeline RAG documentaire : chunking, embeddings, recherche sémantique et base vectorielle ChromaDB.}
        \item {Développement d’un outil de collecte automatisée de métadonnées d’articles scientifiques à partir de sources publiques.}
        \item {Mise en place de webhooks pour transmettre les résultats de tâches planifiées (cron).}
        \item {Conteneurisation de certains outils avec Docker et mise en place sur GitHub d’un pipeline CI/CD automatisant l’intégration.}
        \item {Comparaison de Qwen, Gemma et GLM selon leur latence, leur TTFT et la qualité de leurs réponses. Rédaction d’un article scientifique sur le protocole et les résultats.}
      \end{cvitems}
    }

  \cventry
    {Stagiaire de recherche en apprentissage par renforcement}
    {Université Karunya}
    {Coimbatore, Tamil Nadu, Inde}
    {Juin 2025 -- Sept. 2025}
    {
      \begin{cvitems}
        \item {Conception, développement et comparaison d’une planification de trajectoire déterministe et d’une approche fondée sur l’apprentissage par renforcement pour la navigation autonome d’un robot hospitalier.}
        \item {Déploiement et optimisation du logiciel et du modèle d’apprentissage par renforcement sur un Raspberry Pi embarqué.}
      \end{cvitems}
    }
\end{cventries}
